% ============================================================================
% AERIS Paper - Section 4: AERIS Protocol Design
% Target Journal: Applied Intelligence (APIN) - Springer
% ============================================================================

\section{AERIS Protocol Design}
\label{sec:protocol}

This section presents the AERIS protocol architecture, detailing its
three-layer design that achieves $O(\log n)$ transmission latency while
maintaining high reliability and lightweight operation.

\subsection{Protocol Overview}

AERIS operates through a hierarchical three-layer architecture, as
illustrated in Figure~\ref{fig:architecture}:

\begin{enumerate}
    \item \textbf{Layer 1: Context-Adaptive Switching (CAS)}: Selects
    optimal transmission mode for intra-cluster communication.

    \item \textbf{Layer 2: Skeleton Backbone}: Establishes stable
    inter-cluster paths using PCA-based principal axis selection.

    \item \textbf{Layer 3: Gateway Coordination}: Manages cluster-head
    to base-station transmission with load balancing.
\end{enumerate}

The key insight enabling $O(\log n)$ latency is the \textbf{parallel
hierarchical transmission}: while PEGASIS requires sequential chain
traversal ($O(n)$ hops), AERIS parallelizes transmission at each
hierarchical level. Data flows from nodes to cluster heads (parallel),
through the skeleton backbone ($O(\log n)$ hops), to gateways, and
finally to the base station.

\begin{figure}[htbp]
\centering
\includegraphics[width=0.95\columnwidth]{fig_aeris_architecture}
\caption{AERIS Three-Layer Architecture. Data flows upward through
parallel intra-cluster aggregation (Layer 1), skeleton backbone routing
(Layer 2), and gateway-to-BS transmission (Layer 3). Total latency is
$O(\log n)$ compared to PEGASIS's $O(n)$.}
\label{fig:architecture}
\end{figure}

\subsection{Layer 1: Context-Adaptive Switching (CAS)}

The CAS module dynamically selects among three transmission modes based
on real-time cluster conditions:

\begin{itemize}
    \item \textbf{Direct Mode}: Member nodes transmit directly to cluster
    head. Optimal for small clusters with good link quality.

    \item \textbf{Chain Mode}: Data is aggregated along a short chain
    within the cluster. Optimal for elongated cluster geometries.

    \item \textbf{Two-Hop Mode}: Uses an intermediate relay node.
    Optimal when direct link quality is poor.
\end{itemize}

\subsubsection{Mode Selection Algorithm}

For each cluster, CAS computes scores for each mode using a linear
combination of six features:

\begin{equation}
S_{mode} = \sum_{i=1}^{6} w_i \cdot f_i
\label{eq:cas_score}
\end{equation}

where the features $f_i$ and weights $w_i$ are defined in
Table~\ref{tab:cas_features}.

\begin{table}[htbp]
\centering
\caption{CAS Feature Weights by Transmission Mode}
\label{tab:cas_features}
\begin{tabular}{@{}lccc@{}}
\toprule
Feature & Direct & Chain & Two-Hop \\
\midrule
Residual Energy & 0.30 & 0.40 & 0.25 \\
Link Quality & 0.25 & 0.15 & 0.20 \\
Distance to CH & -0.15 & -0.10 & 0.10 \\
Cluster Size & -0.10 & 0.20 & 0.15 \\
Node Density & 0.10 & 0.05 & 0.20 \\
Fairness Index & 0.20 & 0.10 & 0.10 \\
\bottomrule
\end{tabular}
\end{table}

The mode with the highest score is selected:
\begin{equation}
mode^* = \arg\max_{m \in \{direct, chain, twohop\}} S_m
\label{eq:mode_selection}
\end{equation}

\subsubsection{Complexity Analysis}

CAS mode selection requires computing three scores, each involving 6
multiplications and 5 additions:

\begin{equation}
T_{CAS} = O(1) \quad \text{(constant time per cluster)}
\label{eq:cas_complexity}
\end{equation}

Total CAS overhead for $k$ clusters: $O(k)$, where $k \ll n$ typically.

\subsection{Layer 2: Skeleton Backbone Network}

The skeleton backbone establishes efficient inter-cluster routing paths
using PCA-based geometric analysis.

\subsubsection{Principal Axis Selection}

Given $k$ cluster head positions $\{(x_1, y_1), \ldots, (x_k, y_k)\}$,
we compute the principal axis of the network using PCA:

\begin{enumerate}
    \item Compute centroid: $(\bar{x}, \bar{y}) = \frac{1}{k}\sum_{i=1}^{k}(x_i, y_i)$

    \item Compute covariance matrix:
    \begin{equation}
    \mathbf{C} = \frac{1}{k}\sum_{i=1}^{k}
    \begin{pmatrix}
    (x_i - \bar{x})^2 & (x_i - \bar{x})(y_i - \bar{y}) \\
    (x_i - \bar{x})(y_i - \bar{y}) & (y_i - \bar{y})^2
    \end{pmatrix}
    \label{eq:covariance}
    \end{equation}

    \item Extract principal eigenvector $\mathbf{v}_1$ (direction of
    maximum variance)

    \item Project cluster heads onto principal axis and select $k_{skel}$
    nodes closest to the axis as skeleton nodes
\end{enumerate}

\subsubsection{Skeleton Node Selection}

Skeleton nodes are selected based on a composite score:

\begin{equation}
G_{skel}(i) = \alpha \cdot prox_i + \beta \cdot cent_i + \gamma \cdot E_i
\label{eq:skeleton_score}
\end{equation}

where:
\begin{itemize}
    \item $prox_i$: Proximity to principal axis (normalized)
    \item $cent_i$: Network centrality (closeness centrality)
    \item $E_i$: Normalized residual energy
    \item $\alpha = 0.4$, $\beta = 0.35$, $\gamma = 0.25$ (default weights)
\end{itemize}

\subsubsection{Routing Through Skeleton}

Non-skeleton cluster heads route to the nearest skeleton node, which
then forwards data along the backbone toward the gateway. The skeleton
forms a tree structure with depth $O(\log k)$, where $k$ is the number
of cluster heads.

\subsubsection{Complexity Analysis}

\begin{itemize}
    \item PCA decomposition: $O(k^2)$ for $2 \times 2$ covariance matrix
    \item Axis proximity scoring: $O(k)$
    \item Centrality computation: $O(k^2)$ using Floyd-Warshall
    \item Top-$k_{skel}$ selection: $O(k \log k)$
\end{itemize}

Total skeleton setup complexity:
\begin{equation}
T_{skeleton} = O(k^2)
\label{eq:skeleton_complexity}
\end{equation}

Since $k = O(\sqrt{n})$ typically (with 10\% cluster head probability),
we have $T_{skeleton} = O(n)$.

\subsection{Layer 3: Gateway Coordination}

The gateway layer manages the final hop from cluster heads to the base
station, incorporating load balancing to prevent hotspots.

\subsubsection{Gateway Selection}

Gateways are selected from skeleton nodes based on:

\begin{equation}
G_{gw}(i) = -\lambda_1 \cdot d_{BS}(i) + \lambda_2 \cdot cent_i + \lambda_3 \cdot E_i - \lambda_4 \cdot load_i
\label{eq:gateway_score}
\end{equation}

where:
\begin{itemize}
    \item $d_{BS}(i)$: Distance to base station
    \item $cent_i$: Network centrality
    \item $E_i$: Residual energy
    \item $load_i$: Current traffic load (for load balancing)
    \item $\lambda_1 = 0.35$, $\lambda_2 = 0.25$, $\lambda_3 = 0.25$,
    $\lambda_4 = 0.15$
\end{itemize}

\subsubsection{Load Limiting Mechanism}

To prevent gateway overload, AERIS enforces:

\begin{equation}
load_i \leq L_{max} \quad \forall i \in Gateways
\label{eq:load_limit}
\end{equation}

where $L_{max}$ is the configurable \texttt{gateway\_load\_limit}
parameter. When a gateway reaches capacity, traffic is redirected to
the next-best gateway, ensuring balanced energy consumption.

\subsubsection{Complexity Analysis}

Gateway selection: $O(k_{skel})$ for scoring and selection, where
$k_{skel} \ll k$.

\subsection{Complete Protocol Flow}

Algorithm~\ref{alg:aeris} presents the complete AERIS protocol flow
for one communication round.

\begin{algorithm}[htbp]
\caption{AERIS Protocol - One Round}
\label{alg:aeris}
\begin{algorithmic}[1]
\REQUIRE Network $G = (V, E)$, Base station $BS$
\ENSURE Data delivered to $BS$

\STATE \textbf{// Phase 1: Cluster Formation}
\STATE $CH \leftarrow$ SelectClusterHeads($V$, $p_{CH}$)
\STATE AssignMembers($V$, $CH$)

\STATE \textbf{// Phase 2: Skeleton Setup}
\STATE $axis \leftarrow$ ComputePrincipalAxis($CH$)
\STATE $Skel \leftarrow$ SelectSkeletonNodes($CH$, $axis$, $k_{skel}$)

\STATE \textbf{// Phase 3: Gateway Selection}
\STATE $GW \leftarrow$ SelectGateways($Skel$, $BS$, $k_{gw}$)

\STATE \textbf{// Phase 4: Data Collection (Parallel)}
\FOR{each cluster $c$ \textbf{in parallel}}
    \STATE $mode \leftarrow$ CAS\_SelectMode($c$)
    \STATE CollectData($c$, $mode$)
\ENDFOR

\STATE \textbf{// Phase 5: Backbone Routing}
\FOR{each non-skeleton CH $ch$}
    \STATE RouteToNearestSkeleton($ch$, $Skel$)
\ENDFOR

\STATE \textbf{// Phase 6: Gateway Transmission}
\FOR{each skeleton node $s$}
    \STATE $gw \leftarrow$ SelectBestGateway($s$, $GW$, $L_{max}$)
    \STATE RouteToGateway($s$, $gw$)
\ENDFOR

\STATE \textbf{// Phase 7: Base Station Delivery}
\FOR{each gateway $gw$ \textbf{in parallel}}
    \STATE TransmitToBS($gw$, $BS$)
\ENDFOR

\end{algorithmic}
\end{algorithm}

\subsection{Latency Analysis}

\begin{theorem}[AERIS Latency Bound]
For a network with $n$ nodes and $k$ cluster heads, AERIS achieves
end-to-end transmission latency of $O(\log n)$.
\end{theorem}

\begin{proof}
The transmission path consists of:
\begin{enumerate}
    \item Intra-cluster: 1--2 hops (CAS mode dependent) $= O(1)$
    \item CH to skeleton: 1 hop $= O(1)$
    \item Skeleton backbone: $O(\log k)$ hops (tree depth)
    \item Gateway to BS: 1 hop $= O(1)$
\end{enumerate}

With $k = O(\sqrt{n})$ cluster heads, the skeleton depth is
$O(\log \sqrt{n}) = O(\log n)$.

Total latency: $O(1) + O(1) + O(\log n) + O(1) = O(\log n)$. \qed
\end{proof}

\subsection{Comparison with PEGASIS}

Table~\ref{tab:complexity_comparison} compares AERIS and PEGASIS
complexity characteristics.

\begin{table}[htbp]
\centering
\caption{Complexity Comparison: AERIS vs PEGASIS}
\label{tab:complexity_comparison}
\begin{tabular}{@{}lcc@{}}
\toprule
Metric & AERIS & PEGASIS \\
\midrule
Transmission Latency & $O(\log n)$ & $O(n)$ \\
Setup Complexity & $O(n)$ & $O(n^2)$ \\
Per-Round Overhead & $O(k)$ & $O(n)$ \\
Memory Requirement & $O(k)$ & $O(n)$ \\
Chain Reconstruction & Not needed & $O(n^2)$ per failure \\
\bottomrule
\end{tabular}
\end{table}

At 500 nodes, this translates to:
\begin{itemize}
    \item AERIS latency: $\sim 11$ hops $\approx$ 110ms
    \item PEGASIS latency: $\sim 250$ hops $\approx$ 2500ms
    \item \textbf{AERIS achieves 96\% latency reduction}
\end{itemize}
